\section{Direct strict Selmer vanishing over the original cubic base}
\label{sec:cubic-original-strict-vanishing}

The $\delta=1$ part of the certified calculation in AC-G103 has a
stronger consequence than its four-sector dimension estimate.
It makes the strict residual group over the original field $K$ zero.
An elementary coefficient argument then makes the full discrete
strict group zero, over its original coefficient DVR.  Neither step
uses an automorphic endpoint, an auxiliary host, a deformation model,
or a mixed Bloch--Kato vanishing premise.

\paragraph{Conventions and the exact local conditions.}
Put $p=269$, $K=\mathbb Q(\sqrt{-3})$, and let
$\chi:G_K\to\mu_3$ be either fixed generator character of the
cyclic extension $K(\sqrt[3]{269})/K$.
Write $c_p:G_{\mathbb Q}\to\mathbb Z_p^\times$ for the actual
action on $p$-power roots of unity, and put $\psi=\chi_{-4}$.
For an integer $e\ge1$ satisfying $e\equiv267\pmod{268}$ and
$j\in\{1,2\}$, define
\[
 a_j=\chi^{-j}\psi^{-1}c_p^{-e}|_{G_K}.
 \tag{G114.1}
\]
Let $E_p/\mathbb Q_p$ be any finite coefficient extension containing
the values of $\chi$, with valuation ring $O$, uniformizer $\varpi$,
and residue field $k$.  The character in (G114.1) takes values in
$O^\times$.  Put $D_j=(E_p/O)(a_j)$ with its discrete topology and
$\bar a_j:G_K\to k^\times$ for the reduction of $a_j$.
For either of these coefficient modules $V$ put
\[
 \begin{split}
 \mathcal H_p(V)&=\prod_{v\mid p}H^1_{\mathrm{cts}}(G_{K_v},V),\\
 \mathcal H_{\ne p}(V)&=\prod_{v\nmid p,\infty}
                                 H^1_{\mathrm{cts}}(I_v,V),
 \end{split}
\]
and use precisely
\[
 \operatorname{Sel}_{\rm str}(K,V)=\ker\bigl(
 H^1_{\mathrm{cts}}(G_K,V)\longrightarrow
                    \mathcal H_p(V)\times\mathcal H_{\ne p}(V)\bigr).
 \tag{G114.2}
\]
Thus the coefficient-prime condition is zero in the full local
cohomology group; away from $p$ it is the unramified condition.
This is the original-field strict group in AC-G113.
No Bloch--Kato finite condition is substituted for its local zero
condition.  At the complex places of $K$ there is no further condition.

\begin{theorem}[AC-G114-residual: the original strict residual group is zero]
\label{thm:ACG114residual}
For every datum in (G114.1), every finite field $k$ of characteristic
$269$ containing its character values, and both $j=1,2$,
\[
          \operatorname{Sel}_{\rm str}(K,k(\bar a_j))=0.
 \tag{G114.3}
\]
The only computational input is the already certified class bound
for $\mathbb Q(\sqrt[3]{269},\sqrt3)$ and the exact single-unit
calculation in the $D=3$, $\delta=1$ sector of AC-G103.
The starting-endpoint and host premises in the automorphic part of
AC-G103-certificate are not premises of (G114.3).
\end{theorem}

\begin{proof}
\emph{1. Identify the coefficient and its exact Kummer field.}
Since $-e\equiv1\pmod{268}$ and $\psi^{-1}=\psi$,
\[
       \bar a_j=\bar\chi^{-j}\psi\bar c_p,
       \eta_j:=\bar a_j/\bar c_p=\bar\chi^{-j}\psi.
 \tag{G114.4}
\]
Put $t=\sqrt[3]{269}$, $s=\sqrt3$, and
\[
 E_c=K(t),\qquad M=K(t,s)=K(t,i),\qquad L=\mathbb Q(t,s).
 \tag{G114.5}
\]
The equality follows from $\sqrt{-3}=is\in K$.
The field cut out over $K$ by $\eta_j$ is precisely $M$:
the cubic and quadratic factors have coprime orders and cut out
$E_c/K$ and $K(i)/K$, respectively.
Consequently $\Delta=\operatorname{Gal}(M/K)=C_3\times C_2$.
The nontrivial element $b$ of $C_2$ fixes $E_c$ and sends $s$ to
$-s$; in particular $\eta_j(b)=-1$.
This is exactly $\delta=1$ and $D=3$ in (G103C.11), rather than
the untwisted $D=1$ cubic field.

Let $S_M$ be the primes of $M$ above $p$ and write
\[
 \begin{split}
 V_{M,S_M}&=\{x\in M^\times:
      p\mid\operatorname{ord}_{\mathfrak q}(x)
                   \text{ for all }\mathfrak q\nmid p\}/M^{\times p},\\
 \mathcal L_M&=\bigoplus_{w\mid p}M_w^\times/M_w^{\times p}.
 \end{split}
 \tag{G114.6}
\]
The pure Kummer descent of (G103.4), or directly prime-to-$p$
inflation--restriction and the Kummer sequence, gives
\[
 \operatorname{Sel}_{\rm str}(K,k(\bar a_j))
  \simeq e_{\eta_j^{-1}}\ker\left(
       V_{M,S_M}\otimes_{\mathbb F_p}k
            \longrightarrow\mathcal L_M\otimes_{\mathbb F_p}k\right).
 \tag{G114.7}
\]
Here the natural action on $H^1(M,\mu_p)$ is twisted by $\eta_j$,
so invariance selects the \emph{inverse} character
$\eta_j^{-1}=\bar\chi^j\psi$ in the Kummer factor.
Restriction to a local decomposition group, and to inertia, has
prime-to-$p$ quotient and therefore preserves the zero-class
conditions.  Away from $p$, Kummer unramifiedness is valuation
divisible by $p$; at $p$, zero means an actual local $p$th power.
No adjunction of $\mu_p$ to $M$ is required.

\emph{2. The target class part is zero.}
The retained default-flag PARI 2.19.1 certificate is for
\[
 P(Y)=Y^6-9Y^4+538Y^3+27Y^2+4842Y+72334,
 \quad Y=s-t,\quad h_L=7.
 \tag{G114.8}
\]
It has recorded output \texttt{CERTIFY=1} followed by
\texttt{DONE}.  This is the unconditional full class-group
certificate already used in (G103C.13), not a provisional
\texttt{bnfinit} output or the separate $D=115$ flag-one assertion.
The field identification is exact: expansion gives $P(s-t)=0$ and
\[
       s=\frac{Y^3+9Y+269}{3Y^2+3}.
\]
The denominator is nonzero at the real embedding and $[L:\mathbb Q]=6$.

For clarity, the target class deduction can be made without a
class-number product.  One has
$\operatorname{Gal}(M/\mathbb Q)=S_3\times\langle b\rangle$.
Let $r$ be the reflection in $S_3$ fixing $t$; it also fixes $s$,
so $M^{\langle r\rangle}=L$.
Extension and norm on ideal classes give
\[
       \operatorname{Cl}(M)[p]^{\langle r\rangle}
                   \simeq\operatorname{Cl}(L)[p]=0,
 \tag{G114.9}
\]
also after extension of scalars to $k$; the degree two is invertible.
If a nonzero class $x$ had character $\eta_j^{-1}$ for $\Delta$,
then $rx$ would have the other nontrivial cubic character and the
same $b$-sign.  The distinct cubic eigenspaces imply $x+rx\ne0$.
But $x+rx$ is $r$-invariant, contradicting (G114.9).
Thus the target character part of $\operatorname{Cl}(M)[p]$ is zero.

The same conclusion holds for the $S_M$-class group.  One may use
the exact integral character-projector argument in G103: the
kernel of ordinary-to-$S_M$ classes is generated by $p$-primes,
and the target projector kills their permutation module because
the full cubic group is inertia at $p$.
Equivalently, apply the prime-to-$p$ idempotent to the finite
$p$-primary class group first; a nonzero finite target component
would have nonzero $p$-torsion, which was just excluded, and its
$S_M$ quotient has zero target component as well.
The Kummer unit/class sequence therefore identifies the target
part of $V_{M,S_M}$ with its $S_M$-unit part.

\emph{3. The one-dimensional unit part has injective localization.}
The unit multiplicity calculation of G103 is also independent
of every host and endpoint.  Since $K$ is imaginary quadratic,
the complex places of $M$ form the regular $\Delta$-set.
The logarithmic unit theorem gives the regular representation
minus the trivial character.  There is no $269$-torsion unit,
since $[M:\mathbb Q]=12<268$.
Prime-to-$p$ integral character idempotents preserve these
multiplicities modulo $p$, so every nontrivial character occurs
once.  The additional $S_M$-unit permutation term has no target
part because cubic inertia at $p$ is nontrivial on that character.
The $b=-1$, nontrivial-cubic part is consequently two-dimensional
over $\mathbb F_{269}$ and has one line of each target character
over $k$.  It is simple over $\mathbb F_{269}[C_3]$, since
$X^2+X+1$ is irreducible modulo $269$.

Here is the exact nonzero local image retained from (G103C.15)--(20).
In $L=\mathbb Q(t,s)$ the element
\[
\begin{split}
 U={}&\frac{1437754535822440}{3}+209434299120299s
 -\frac{172020383324665}{3}t+108412322108067ts\\
 &+\frac{83072025828619}{3}t^2-4321330156673t^2s
\end{split}
 \tag{G114.10}
\]
has monic integral multiplication polynomial
\[
 Z^6-AZ^5+BZ^4-CZ^3+BZ^2-AZ+1,
 \tag{G114.11}
\]
where $A=2875509071644880$,
$B=7881560119642923513355187503223$, and
$C=17678805150002035094958926455485206906030719856$.
Exact rational matrix multiplication verifies this polynomial;
its constant term one makes both $U$ and $U^{-1}$ integral.
Set
\[
             v=U/b(U),\qquad
             w=v^3/N_{M/K(s)}(v).
 \tag{G114.12}
\]
Then $b(w)=w^{-1}$ and $N_{C_3}(w)=1$, so its mod-$p$ class
belongs to the two-dimensional unit part just described.
At a prime above $p$, normalize $v(t)=1$, hence $v(p)=3$.
The residue field is $\mathbb F_{p^2}$, with $s^2=3$ nonsquare
over $\mathbb F_p$.  Direct reduction of (G114.10) gives
\[
 U=A_0+B_0t\pmod{t^2},\quad
 A_0=82+122s,\quad B_0=260+95s,\quad B_0/A_0=11s.
 \tag{G114.13}
\]
Thus $v$ has logarithmic linear coefficient $22s$.
Its cubic norm has no $t$ term and has constant term equal to
the cube of the constant term of $v$.  Consequently
\[
                         w=1+66s\,t\pmod{t^2}.
 \tag{G114.14}
\]
Because the absolute ramification index is three and $3<p-1$,
logarithm and exponential identify $1+\mathfrak m$ with
$\mathfrak m$; multiplication by $p$ has image
$\mathfrak m^4$.  The residue-unit group has order $p^2-1$,
prime to $p$.  Equation (G114.14) therefore shows that $w$ is
not a local $p$th power.

Localization into the \emph{sum of all} $p$-completions is
$\Delta$-equivariant.  It is nonzero on the simple
two-dimensional unit part by (G114.14), hence injective there.
After extending scalars to $k$ it is injective on both target
cubic lines.  This also handles the possible exchange of two
$p$-primes by $b$; no individual completion is falsely declared
$\Delta$-stable.  Combining the class calculation with (G114.7)
proves (G114.3) for both $j$.
All these computations are over $K$, $M$, and $L$ and require
none of the other three quadratic host sectors of G103.
\end{proof}

\begin{theorem}[AC-G114-strict: vanishing with the original discrete coefficients]
\label{thm:ACG114strict}
For every finite coefficient DVR $O$ and all data in (G114.1),
\[
                    \operatorname{Sel}_{\rm str}(K,D_j)=0
                       \qquad(j=1,2).
 \tag{G114.15}
\]
In particular its $O$-length and its normalized length
$\operatorname{length}_O/e_O$, where
$e_O=\operatorname{length}_O(O/pO)$, are zero.
The conclusion needs no condition $e_O=1$ and no extension or
replacement of the original coefficient ring.
\end{theorem}

\begin{proof}
\emph{1. Local invariants and the uniformizer kernel.}
Multiplication by $\varpi$ on $D_j$ is surjective, and
\[
 D_j[\varpi]\simeq k(\bar a_j),\qquad
 \bar u\longmapsto u/\varpi+O.
 \tag{G114.16}
\]
This uses a uniformizer of $O$, not multiplication by $p$;
thus (G114.16) remains valid when $O/\mathbb Z_p$ is ramified.

There is a unique prime of $K$ above $p$.
The cubic extension $K_p(t)/K_p$ is totally ramified of degree
three, so $\bar\chi$ has inertia image of order three.
The rational cyclotomic reduction has inertia order dividing
$p-1=268$, whereas $\psi$ is unramified at $p$.
The coprime orders imply that $\bar a_j$ is nontrivial on
coefficient-prime inertia.  Explicitly, raise an inertia element
with nontrivial cubic value to the $268$th power: its
$\bar c_p$ value becomes one and its cubic value remains
nontrivial, since $268\equiv1\pmod3$.
For such an element $g$, $a_j(g)-1$ is a unit in $O$.
Multiplication by this unit on $E_p/O$ is invertible, so
\[
                     D_j^{G_{K_p}}=0,
                     \qquad D_j^{G_K}=0.
 \tag{G114.17}
\]
In particular both invariant modules are $\varpi$-divisible.

At a finite place $v\nmid p$, the full $p$-adic cyclotomic
character is trivial on inertia.  The remaining character
$\chi^{-j}\psi^{-1}$ on $I_v$ has finite order dividing six,
which is prime to $p$.
If it is trivial, $D_j^{I_v}=D_j$ is $\varpi$-divisible.
If it is nontrivial, some value is a nontrivial root of unity
of order prime to $p$.  Its difference from one is a unit in
$O$, so $D_j^{I_v}=0$.
Thus
\[
             D_j^{I_v}/\varpi D_j^{I_v}=0
                         \quad\text{for every }v\nmid p.
 \tag{G114.18}
\]
This checks all away places, including those ramified in
$\psi$ or $\chi$, without omitting any inertia condition.

\emph{2. Lift a uniformizer-torsion strict class to the residual group.}
The continuous cohomology sequence of
$0\to D_j[\varpi]\to D_j\xrightarrow{\varpi}D_j\to0$
gives
\[
       H^1(G_K,D_j[\varpi])\simeq H^1(G_K,D_j)[\varpi],
 \tag{G114.19}
\]
where (G114.17) gives injectivity.
For every local group in (G114.2), (G114.17)--(18) give
injectivity of
$H^1(G_{K_p},D_j[\varpi])\to H^1(G_{K_p},D_j)$ and
$H^1(I_v,D_j[\varpi])\to H^1(I_v,D_j)$.
Hence the unique lift in (G114.19) of a strict class is itself
strict.  Equations (G114.3) and (G114.16) imply
\[
                \operatorname{Sel}_{\rm str}(K,D_j)[\varpi]=0.
 \tag{G114.20}
\]

Here is a direct cocycle verification of the exactness and the
local step, to expose every coefficient hypothesis.
If $\varpi[z]=0$, choose $m\in D_j$ with
$\varpi z=dm$, and choose $n\in D_j$ with $\varpi n=m$.
Then $z-dn$ takes values in $D_j[\varpi]$.
If such a residual cocycle is a coboundary $du$ in $D_j$,
then $\varpi u$ is invariant.  For the global group that
invariant is zero by (G114.17), so $u\in D_j[\varpi]$.
For a local group, divisibility of its invariants provides an
invariant $u_0$ with $\varpi u_0=\varpi u$.
Then $u-u_0\in D_j[\varpi]$ and $du=d(u-u_0)$.
This proves the required local injectivity and preserves the
actual zero restrictions.  All the displayed coboundaries are
continuous, as the coefficient action is continuous and discrete.

\emph{3. Every continuous class is uniformizer-primary torsion.}
A continuous cocycle from the compact profinite group $G_K$
to the discrete module $D_j$ has finite image.
Each element of $E_p/O$ is killed by a power of $\varpi$;
one power therefore kills all values of that cocycle.
Thus every element of $H^1(G_K,D_j)$, and of its strict
submodule, is killed by a power of $\varpi$.
If a nonzero strict class existed, choose its least annihilating
power $\varpi^n$.  Its multiple by $\varpi^{n-1}$ would be a
nonzero element of the group in (G114.20), a contradiction.
This proves (G114.15).  No finiteness, cotorsion, Iwasawa
invariant, automorphy, or deformation-rigidity hypothesis has
entered this argument.
\end{proof}

\paragraph{Exact consequence and boundary.}
On $N=4$, $\psi=\chi_{-4}$, $e\equiv267\pmod{268}$,
the strict arithmetic quantity $s_j$ in AC-G113 is therefore
zero directly, for each original coefficient ring and both
cubic characters.  Any stricter subgroup obtained by adding
further local zero conditions is also zero.
The conditional model-and-counting theorem AC-G113 remains
correct, but its premises are unnecessary for this particular
strict vanishing conclusion.
This does not prove M4, vanishing of the mixed characteristic-zero
Bloch--Kato group, an automorphic endpoint's existence, or the
same result for other characters, other weight residue classes,
other coefficient primes, or relaxed local conditions.

\paragraph{Certificate, source and geometric scope.}
The class input is the already certified $D=3$ output, its exact
defining polynomial, and the literal default-flag
\texttt{bnfcertify} scope of the official PARI 2.19.1 manual.
The single-unit polynomial and finite-field calculation are
independently reproducible by the retained exact-arithmetic script;
they do not assume a computed unit basis is complete.
The Kummer input is Stacks Project, Section 59.28, Lemma 59.28.1
(tag 03PK), its full proof, and the ensuing unit/Picard sequence.
The logarithmic unit-representation and prime-to-$p$ descent
arguments are the previously checked pure arithmetic parts of
AC-G103, rederived above in their exact $D=3$ scope.

Geometrically, the residual kernel classifies the
$\eta_j^{-1}$-character part of Kummer torsors on
$\operatorname{Spec}\mathcal O_M[1/p]$, with trivial class over
every $p$-adic completion.  The certified class part and the
injective unit localization show that this part has no nontrivial
torsor class.  The passage to $E_p/O$ is the explicit
uniformizer-kernel argument just proved.  This description adds
no new motive, regulator nonvanishing, or Hecke-support assertion.
